\ifdefined\iscombined
	\lecturetitle{Faults}
\else
\documentclass{article}
\usepackage{listings}
\usepackage{amsthm}
\usepackage{comment}
\usepackage{amsmath}
\usepackage{amsfonts}
\usepackage{sectsty}
\usepackage{hyperref}
\usepackage{fancyhdr}
\usepackage[top=1in,bottom=1in,left=1in,right=1in]{geometry}
\usepackage{float}
\usepackage{graphicx}
\usepackage{tabularx}
\usepackage{mathtools}

\date{
	\vspace{-.75em}
	{\large \today}
}

\title{
	\vspace*{-1.5em}
	{\Huge Lecture 2} \\
	\vspace{-1em}
	\rule{\textwidth/4}{.01em} \\
	\vspace{-.25em}
	{\Large Faults}
	\vspace{-.75em}
}

\author{
	{\large Matthew Farah}
}

\hypersetup{
	colorlinks=true,
	linkcolor=blue,
	filecolor=magenta,
	urlcolor=blue,
	citecolor=blue,
	anchorcolor=blue
}

\fancyhead{}
\fancyhead[L]{\today}
\fancyhead[R]{Lecture 2 - Faults}

\sectionfont{\fontsize{12}{12}\bfseries}
\subsectionfont{\fontsize{10}{12}\normalfont}
\subsubsectionfont{\fontsize{10}{12}\normalfont}

\begin{document}

\maketitle
\pagestyle{fancy}

\fi

\begin{itemize}
	\item Testing is defined as the process of evaluating software to estimate how well it meets some set of criteria.
	\begin{itemize}
		\item One of the primary goals of testing is to identify `faults' in a program.
	\end{itemize}
	\item Faults in programs are defined as anything which may cause the program to behave in an unintended way during execution.
	\begin{itemize}
		\item Synonymous with `bug' or `defect'.
	\end{itemize}
	\item `Errors' and `failures' are defined as specific kinds of faults.
	\begin{itemize}
		\item A program is said to error if it enters a bad state during its execution.
		\item A failure is an observed instance of the program deviating from expected behaviour during execution.
	\end{itemize}
	\item Tests can only be used to prove or disprove the existence of faults in code.
	\item A test case is a specific choice of inputs with a given set of initial conditions designed to verify if the program works as intended.
	\begin{itemize}
		\item Test cases use `test oracles' to determine this.
		\begin{itemize}
			\item Test oracles compare the actual and expected result of the program in a test case.
		\end{itemize}
	\end{itemize}
	\item Edge case inputs are synonymous with boundary values.
	\item Tests should be designed systematically, using domain knowledge about the program's inputs.
\end{itemize}

\ifdefined\iscombined
	\newpage
\else
	\end{document}
\fi
